\begin{question}{Representation Surgery}

\begin{subquestion}

\subsubsection*{Linear Guarding vs.\ Arbitrary Classifiers}

No. By Definition 2, the statement that $\mathbf{X}$ \emph{linearly guards} $Z$ only guarantees that the trivially attainable loss $L_\tau$ is optimal over the set of all \textbf{affine transformations} (linear predictors). 
Mathematically, it ensures that $\expectation\left[\mathcal{L}(f(\mathbf{X}), Z)\right] \geq L_\tau$ holds for any affine predictor $f(\mathbf{x}) = \mathbf{A}\mathbf{x} + \mathbf{b}$.
However, the classifier $c : \mathbf{x} \mapsto c(\mathbf{x}) \in \mathcal{Z}$ proposed in the question is an \emph{arbitrary} classifier, meaning it is not restricted to being linear.
Because linear guarding only constrains linear probes, a more complex, non-linear classifier such as a multi-layer perceptron could potentially draw non-linear decision boundaries to extract the concept information $Z$ still embedded within $\mathbf{X}$.
Consequently, such a non-linear classifier could successfully achieve a loss lower than the trivially attainable loss $L_\tau$.

\end{subquestion}

\begin{subquestion}

\subsubsection*{Class-Conditional Mean Equality Implies Linear Guarding}

If each class-conditional mean satisfies $\expectation[\mathbf{X} \mid Z = i] = \expectation[\mathbf{X}]$, then $\mathbf{X}$ linearly guards $Z$ w.r.t. the cross-entropy loss $\mathcal{L}$. 
To prove that $\mathbf{X}$ linearly guards $Z$, we must show that for any affine predictor $f(\mathbf{x}) = \mathbf{A}\mathbf{x} + \mathbf{b}$ (where $\mathbf{A} \in \R^{k \times d}$ and $\mathbf{b} \in \R^k$), the expected loss is bounded by the trivially attainable loss $L_\tau$:
\begin{equation}
    \expectation_{\mathbf{X}, Z}\left[\mathcal{L}(\mathbf{A}\mathbf{X} + \mathbf{b}, Z)\right] \geq L_\tau
\end{equation}

We can rewrite the expected loss by conditioning on the class $Z$:
\begin{equation}
    \expectation_{\mathbf{X}, Z}\left[\mathcal{L}(\mathbf{A}\mathbf{X} + \mathbf{b}, Z)\right] = \expectation_Z \Big[ \expectation_{\mathbf{X} \mid Z}\left[\mathcal{L}(\mathbf{A}\mathbf{X} + \mathbf{b}, Z) \mid Z\right] \Big]
\end{equation}

The cross-entropy loss $\mathcal{L}(\mathbf{y}, Z = i) = -\ln \frac{\exp(y_i)}{\sum_{j=1}^k \exp(y_j)}$ is convex w.r.t. its input logits $\mathbf{y}$. 
Because $\mathbf{y} = \mathbf{A}\mathbf{x} + \mathbf{b}$ is an affine, the composition $\mathcal{L}(\mathbf{A}\mathbf{x} + \mathbf{b}, Z=i)$ remains convex w.r.t. $\mathbf{x}$.
Applying Jensen's inequality gives:
\begin{equation}
    \expectation_{\mathbf{X} \mid Z=i}\left[\mathcal{L}(\mathbf{A}\mathbf{X} + \mathbf{b}, i)\right] \geq \mathcal{L}\Big(\expectation_{\mathbf{X} \mid Z=i}\left[\mathbf{A}\mathbf{X} + \mathbf{b}\right], i\Big) = \mathcal{L}\Big(\mathbf{A}\,\expectation[\mathbf{X} \mid Z=i] + \mathbf{b}, i\Big)
\end{equation}

Use the fact that class-conditional means equal the global mean, $\expectation[\mathbf{X} \mid Z = i] = \expectation[\mathbf{X}]$.
\begin{equation}
    \expectation_{\mathbf{X} \mid Z=i}\left[\mathcal{L}(\mathbf{A}\mathbf{X} + \mathbf{b}, i)\right] \geq \mathcal{L}\Big(\mathbf{A}\,\expectation[\mathbf{X}] + \mathbf{b}, i\Big)
\end{equation}

Substituting this lower bound back into the outer expectation over $Z$ gives:
\begin{equation}
    \expectation_{\mathbf{X}, Z}\left[\mathcal{L}(\mathbf{A}\mathbf{X} + \mathbf{b}, Z)\right] \geq \expectation_Z \Big[\mathcal{L}(\mathbf{A}\,\expectation[\mathbf{X}] + \mathbf{b}, Z)\Big]
\end{equation}
Notice that $\mathbf{b}' \triangleq \mathbf{A}\,\expectation[\mathbf{X}] + \mathbf{b}$ is simply a constant vector in $\R^k$, independent of $\mathbf{X}$. Thus, the right side of the inequality evaluates the loss of a \emph{constant} predictor $f(\mathbf{x}) = \mathbf{b}'$.

The trivially attainable loss $L_\tau$ is the infimum of the expected loss over all constant predictors:
\begin{equation}
    L_\tau = \inf_{\mathbf{c} \in \R^k} \expectation_Z\left[\mathcal{L}(\mathbf{c}, Z)\right]
\end{equation}
Since $\mathbf{b}'$ is a constant in $\R^k$, its expected loss must be greater than or equal to the infimum:
\begin{equation}
    \expectation_Z \Big[\mathcal{L}(\mathbf{b}', Z)\Big] \geq L_\tau
\end{equation}

Combining the inequalities, we have shown that for any affine predictor $f(\mathbf{x}) = \mathbf{A}\mathbf{x} + \mathbf{b}$:
\begin{equation}
    \expectation_{\mathbf{X}, Z}\left[\mathcal{L}(f(\mathbf{X}), Z)\right] \geq L_\tau
\end{equation}
Therefore, by Definition 2, $\mathbf{X}$ linearly guards $Z$ w.r.t. the cross-entropy loss $\mathcal{L}$.

\end{subquestion}

\begin{subquestion}

\subsubsection*{Boundedness and Symmetry of the Cross-Entropy Gradient}

Let $\mathbf{x}$ be an observed representation, and $\mathbf{y} \triangleq \mathbf{A}\mathbf{x} + \mathbf{b}$.

\paragraph{All partial derivatives of the cross-entropy loss w.r.t. $\mathbf{y}$ are bounded.}

Recall the definition of the cross-entropy loss for logit vector $\mathbf{y}$ and true concept class $Z = c$:
\begin{equation}
    \mathcal{L}(\mathbf{y}, Z = c) = -\ln\left(\frac{\exp(y_c)}{\sum_{j=1}^k \exp(y_j)}\right) = -y_c + \ln\left(\sum_{j=1}^k \exp(y_j)\right)
\end{equation}

We compute the partial derivative of the loss w.r.t. an arbitrary logit $y_m$.
\textbf{Case 1 ($m = c$):}
\begin{align}
    \frac{\partial}{\partial y_c}\mathcal{L}(\mathbf{y}, Z = c)
    &= -1 + \frac{\frac{\partial}{\partial y_c}\sum_{j=1}^k \exp(y_j)}{\sum_{j=1}^k \exp(y_j)} \\
    &= -1 + \frac{\exp(y_c)}{\sum_{j=1}^k \exp(y_j)} = -1 + \mathrm{softmax}(\mathbf{y})_c
\end{align}
Since the softmax output lies in $(0, 1)$ for any $\mathbf{y}$, we have $-1 < \frac{\partial}{\partial y_c}\mathcal{L}(\mathbf{y}, Z = c) < 0$.

\textbf{Case 2 ($m \neq c$):}
\begin{align}
    \frac{\partial}{\partial y_m}\mathcal{L}(\mathbf{y}, Z = c)
    &= 0 + \frac{\frac{\partial}{\partial y_m}\sum_{j=1}^k \exp(y_j)}{\sum_{j=1}^k \exp(y_j)} \\
    &= \frac{\exp(y_m)}{\sum_{j=1}^k \exp(y_j)} = \mathrm{softmax}(\mathbf{y})_m
\end{align}
We have $0 < \frac{\partial}{\partial y_m}\mathcal{L}(\mathbf{y}, Z = c) < 1$. Because partial derivatives fall within the interval $(-1, 1)$, they are uniformly bounded.

\paragraph{For any $j_1, j_2 \neq i$, we have $\dfrac{\partial}{\partial y_i}\mathcal{L}(\mathbf{y}, Z = j_1) = \dfrac{\partial}{\partial y_i}\mathcal{L}(\mathbf{y}, Z = j_2) \neq 0$.}
We find the partial derivative of the loss w.r.t. $y_i$ when the true class is $Z = j$ with $j \neq i$. 
Using the derivation from above, the partial derivatives are the $i$-th component of the softmax output:
\begin{equation}
    \frac{\partial}{\partial y_i}\mathcal{L}(\mathbf{y}, Z = j_1) = \frac{\partial}{\partial y_i}\mathcal{L}(\mathbf{y}, Z = j_2) = \frac{\exp(y_i)}{\sum_{l=1}^k \exp(y_l)}
\end{equation}

Both expressions are equal. Furthermore, because the exponential function $e^x$ is positive for any real $x$, we have $\exp(y_i) > 0$. The denominator is a sum of positive exponentials, so it is greater than zero. Therefore:
\begin{equation}
    \frac{\exp(y_i)}{\sum_{l=1}^k \exp(y_l)} > 0
\end{equation}

Thus, we have shown that:
\begin{equation}
    \frac{\partial}{\partial y_i}\mathcal{L}(\mathbf{y}, Z = j_1) = \frac{\partial}{\partial y_i}\mathcal{L}(\mathbf{y}, Z = j_2) \neq 0
\end{equation}

\end{subquestion}

\begin{subquestion}

\subsubsection*{Linear Guarding Implies Equal Class-Conditional Means}

Because $\mathbf{X}$ linearly guards $Z$, the expected loss for any affine predictor $f(\mathbf{x}) = \mathbf{A}\mathbf{x} + \mathbf{b}$ is bounded by the trivially attainable loss $L_\tau$ achieved by the constant predictor $\mathbf{b}^\ast$.
Thus,
\begin{equation}
    J(\mathbf{A}, \mathbf{b}) \triangleq \expectation_{\mathbf{X}, Z}\left[\mathcal{L}(\mathbf{A}\mathbf{X} + \mathbf{b}, Z)\right]
\end{equation}
attains its minimum at $\mathbf{A} = \mathbf{0}$ and $\mathbf{b} = \mathbf{b}^\ast$.
By the first-order optimality, the gradient of $J$ w.r.t. $\mathbf{A}$, evaluated at $\mathbf{A} = \mathbf{0}, \mathbf{b} = \mathbf{b}^\ast$:
\begin{equation}
    \nabla_{\mathbf{A}}\, \expectation_{\mathbf{X}, Z}\left[\mathcal{L}(\mathbf{A}\mathbf{X} + \mathbf{b}, Z)\right]\Big|_{\mathbf{A} = \mathbf{0}, \mathbf{b} = \mathbf{b}^\ast} = \mathbf{0}
\end{equation}

Move the derivative inside the expectation to evaluate the gradient. Let $\mathbf{y} = \mathbf{A}\mathbf{X} + \mathbf{b}$.
\begin{equation}
    \frac{\partial \mathcal{L}}{\partial A_{ij}} = \frac{\partial \mathcal{L}}{\partial y_i}\, \frac{\partial y_i}{\partial A_{ij}} = \frac{\partial \mathcal{L}}{\partial y_i}\, X_j.
\end{equation}
In part c we proved that the partial derivatives of the cross-entropy loss w.r.t. $\mathbf{y}$ are bounded between $-1$ and $1$.
Because the gradient w.r.t $\mathbf{y}$ is uniformly bounded, the Dominated Convergence Theorem allows us to interchange expectation and differentiation:
\begin{equation}
    \nabla_{\mathbf{A}} J(\mathbf{A}, \mathbf{b}) = \expectation_{\mathbf{X}, Z}\left[\nabla_{\mathbf{A}}\, \mathcal{L}(\mathbf{A}\mathbf{X} + \mathbf{b}, Z)\right] = \expectation_{\mathbf{X}, Z}\left[\nabla_{\mathbf{y}}\, \mathcal{L}(\mathbf{y}, Z)\, \mathbf{X}^\top\right]
\end{equation}

From part c, for any logit $y_m$ the partial derivative of the loss is $\frac{\partial \mathcal{L}}{\partial y_m} = \mathrm{softmax}(\mathbf{y})_m - \mathbb{1}(Z = m)$.
Let $\mathbf{p}(\mathbf{y})$ be the vector of softmax probabilities and let $\widetilde{\mathbf{Z}}$ be the one-hot encoding of $Z$.
\begin{equation}
    \nabla_{\mathbf{y}}\, \mathcal{L}(\mathbf{y}, Z) = \mathbf{p}(\mathbf{y}) - \widetilde{\mathbf{Z}}.
\end{equation}
Evaluating the gradient w.r.t. $\mathbf{A}$ at the optimum $\mathbf{A} = \mathbf{0}, \mathbf{b} = \mathbf{b}^\ast$ so $\mathbf{y} = \mathbf{b}^\ast$ yields:
\begin{equation}
    \expectation_{\mathbf{X}, Z}\left[(\mathbf{p}(\mathbf{b}^\ast) - \widetilde{\mathbf{Z}})\, \mathbf{X}^\top\right] = \mathbf{0}
\end{equation}
Restricting to a single class $i$, this implies
\begin{equation}
    \expectation_{\mathbf{X}, Z}\left[(p_i(\mathbf{b}^\ast) - \mathbb{1}(Z = i))\, \mathbf{X}\right] = \mathbf{0}.
\end{equation}
By linearity of expectation,
\begin{equation}
    p_i(\mathbf{b}^\ast)\, \expectation[\mathbf{X}] = \expectation\left[\mathbf{X}\, \mathbb{1}(Z = i)\right]
\end{equation}

To determine $p_i(\mathbf{b}^\ast)$, recall that $\mathbf{b}^\ast$ minimizes the trivially attainable loss $L_\tau = \expectation_Z[\mathcal{L}(\mathbf{b}, Z)]$. First-order optimality w.r.t. $\mathbf{b}$ gives:
\begin{equation}
    \nabla_{\mathbf{b}}\, \expectation_Z\left[\mathcal{L}(\mathbf{b}, Z)\right] = \expectation_Z\left[\nabla_{\mathbf{b}}\, \mathcal{L}(\mathbf{b}, Z)\right] = \expectation_Z\left[\mathbf{p}(\mathbf{b}) - \widetilde{\mathbf{Z}}\right] = \mathbf{0}
\end{equation}
Hence, at the optimum $\mathbf{b}^\ast$, the softmax probabilities must match the marginal class probabilities:
\begin{equation}
    \mathbf{p}(\mathbf{b}^\ast) = \expectation_Z[\widetilde{\mathbf{Z}}], \qquad\text{and thus}\qquad p_i(\mathbf{b}^\ast) = \mathbb{P}(Z = i).
\end{equation}

Substituting $p_i(\mathbf{b}^\ast) = \mathbb{P}(Z = i)$ back into $(24)$:
\begin{equation}
    \mathbb{P}(Z = i)\, \expectation[\mathbf{X}] = \expectation\left[\mathbf{X}\, \mathbb{1}(Z = i)\right]
\end{equation}
By the definition of conditional expectation, $\expectation[\mathbf{X}\, \mathbb{1}(Z = i)] = \expectation[\mathbf{X} \mid Z = i]\, \mathbb{P}(Z = i)$, so
\begin{equation}
    \mathbb{P}(Z = i)\, \expectation[\mathbf{X}] = \expectation[\mathbf{X} \mid Z = i]\, \mathbb{P}(Z = i).
\end{equation}
Since $\mathbb{P}(Z = i) \neq 0$ by assumption, we can divide both sides by $\mathbb{P}(Z = i)$ to obtain
\begin{equation}
    \expectation[\mathbf{X}] = \expectation[\mathbf{X} \mid Z = i].
\end{equation}
This holds for every class $i$.

\end{subquestion}

\begin{subquestion}

\subsubsection*{Zero Cross-Covariance Characterizes Equal Class-Conditional Means}

The cross-covariance between $\mathbf{X}$ and $\widetilde{\mathbf{Z}}$ is
\begin{equation}
    \boldsymbol{\Sigma}_{\mathbf{X}, \widetilde{\mathbf{Z}}} = \expectation\left[\mathbf{X}\, \widetilde{\mathbf{Z}}^\top\right] - \expectation[\mathbf{X}]\, \expectation[\widetilde{\mathbf{Z}}]^\top.
\end{equation}
Its $j^\text{th}$ column is the covariance of the vector $\mathbf{X}$ with the scalar $\widetilde{Z}_j$:
\begin{equation}
    \mathrm{Cov}(\mathbf{X}, \widetilde{Z}_j) = \expectation[\mathbf{X}\, \widetilde{Z}_j] - \expectation[\mathbf{X}]\, \expectation[\widetilde{Z}_j].
\end{equation}
Since $\widetilde{Z}_j = \mathbb{1}(Z = j)$, we have $\expectation[\widetilde{Z}_j] = \mathbb{P}(Z = j)$ and, by the definition of conditional expectation, $\expectation[\mathbf{X}\, \widetilde{Z}_j] = \expectation[\mathbf{X} \mid Z = j]\, \mathbb{P}(Z = j)$. Substituting gives:
\begin{equation}
    \mathrm{Cov}(\mathbf{X}, \widetilde{Z}_j) = \Big(\expectation[\mathbf{X} \mid Z = j] - \expectation[\mathbf{X}]\Big)\, \mathbb{P}(Z = j)
\end{equation}

\textbf{Forward direction} Assume $\expectation[\mathbf{X} \mid Z = j] = \expectation[\mathbf{X}]$ for every $j$. Then the parenthesised term vanishes, so $\mathrm{Cov}(\mathbf{X}, \widetilde{Z}_j) = \mathbf{0}$ for every column $j$, and therefore $\boldsymbol{\Sigma}_{\mathbf{X}, \widetilde{\mathbf{Z}}} = \mathbf{0}$.

\textbf{Backward direction} Assume $\boldsymbol{\Sigma}_{\mathbf{X}, \widetilde{\mathbf{Z}}} = \mathbf{0}$. Every column must then be zero, so for each $j$:
\begin{equation}
    \Big(\expectation[\mathbf{X} \mid Z = j] - \expectation[\mathbf{X}]\Big)\, \mathbb{P}(Z = j) = \mathbf{0}.
\end{equation}
Since $\mathbb{P}(Z = j) \neq 0$ by assumption, the directions give $\expectation[\mathbf{X} \mid Z = i] = \expectation[\mathbf{X}] \iff \boldsymbol{\Sigma}_{\mathbf{X}, \widetilde{\mathbf{Z}}} = \mathbf{0}$.

\end{subquestion}

\begin{subquestion}

\end{subquestion}

\begin{subquestion}

\end{subquestion}

\begin{subquestion}

\end{subquestion}

\begin{subquestion}

\end{subquestion}

\end{question}